% ===== Chapter 3 opening (~6 min) + section 1: superstition and conspiracy theories (~20 min) =====
\partopener{Գլուխ 3}{Վախ}{Վախենալու բան չկա։ Ինչպես է ուղեղը մեզ անընդհատ վախեցնում}{panic}{peach}

% The one recap frame: webinar 1 already taught fight-or-flight and the amygdala
\begin{frame}{Հիշո՞ւմ ես վեբինար 1-ից}
  \begin{columns}[T]
    \begin{column}{0.31\textwidth}
      \begin{card}[colback=lavender!35!cardwhite]
        \bigemoji[0.8cm]{axe}\par\smallskip
        {\blackfont Խալաթ, թե՞ մարդասպան}\par
        Մթության մեջ ուղեղը նախ վախենում է, հետո՝ ստուգում
      \end{card}
    \end{column}
    \begin{column}{0.31\textwidth}
      \begin{card}[colback=brainpink!30!cardwhite]
        \bigemoji[0.8cm]{bell}\par\smallskip
        {\blackfont Արագ ճանապարհ}\par
        Նշաձև մարմինը (amygdala) ահազանգում է նախքան կեղևը կմտածի
      \end{card}
    \end{column}
    \begin{column}{0.31\textwidth}
      \begin{card}[colback=mint!35!cardwhite]
        \bigemoji[0.8cm]{runner}\par\smallskip
        {\blackfont «Կռվիր կամ փախիր»}\par
        Սիրտը, թոքերը, մկանները՝ մարտական պատրաստության
      \end{card}
    \end{column}
  \end{columns}
  \vspace{3mm}
  \onslide<2->{%
  \begin{card}[colback=peach!45!cardwhite]
    \centering Այսօր՝ ի՞նչ է լինում, երբ այս համակարգը վտանգ է «գտնում» \textbf{ամենուր}
  \end{card}}
\end{frame}

% The world changed, the threat detector did not
\begin{frame}{Աշխարհը փոխվեց, ուղեղը՝ ոչ}
  \begin{columns}[T]
    \begin{column}{0.47\textwidth}
      \begin{card}[colback=mint!35!cardwhite]
        \emoji[1.4em]{tiger}\ {\blackfont Վաղուց}\par\smallskip
        Աշխարհը լի էր իրական, կյանքին սպառնացող վտանգներով։ Ուղեղը հենց դրանք էր որսում։
      \end{card}
    \end{column}
    \begin{column}{0.47\textwidth}
      \onslide<2->{%
      \begin{card}[colback=brainpink!30!cardwhite]
        \emoji[1.4em]{receipt}\ {\blackfont Այսօր}\par\smallskip
        «Վտանգ» կարող է լինել ամեն ինչ. հաշիվը, պետի նամակը, «Ինչու՞ է ոտքս քոր գալիս»։
      \end{card}}
    \end{column}
  \end{columns}
  \vspace{3mm}
  \onslide<3->{%
  \begin{bookquote}
    Այն ամենը, ինչ կարող է բացասական հետևանք ունենալ, որքան էլ փոքր կամ սուբյեկտիվ լինի, գրանցվում է որպես "արժե անհանգստանալ"։
  \end{bookquote}}
\end{frame}

% The menu of the chapter: four ways to fear what is not dangerous
\begin{frame}{Վախի ճաշացանկը}
  \begin{columns}[T]
    \begin{column}{0.48\textwidth}
      \begin{card}[colback=lavender!35!cardwhite]
        \emoji[1.3em]{clover}\ {\blackfont Սնահավատություն}\par
        Ծեսեր, որոնք ոչնչի վրա չեն ազդում, բայց հանգստացնում են
      \end{card}
      \onslide<3->{%
      \begin{card}[colback=mint!35!cardwhite]
        \emoji[1.3em]{spider}\ {\blackfont Ֆոբիաներ}\par
        Գիտես, որ անվնաս է, բայց միևնույնն է՝ սարսափում ես
      \end{card}}
    \end{column}
    \begin{column}{0.48\textwidth}
      \onslide<2->{%
      \begin{card}[colback=butter!45!cardwhite]
        \emoji[1.3em]{ufo}\ {\blackfont Դավադրության տեսություններ}\par
        Տեսականորեն հնարավոր, բայց չափազանց անհավանական վտանգներ
      \end{card}}
      \onslide<4->{%
      \begin{card}[colback=peach!40!cardwhite]
        \emoji[1.3em]{hourglass}\ {\blackfont Անհանգստություն ոչնչից}\par
        «Չափազանց հանգիստ է», «վատ բան պիտի որ լինի»
      \end{card}}
    \end{column}
  \end{columns}
\end{frame}

% When worrying about harmless things starts to harm
\begin{frame}{Երբ անհանգստությունն ինքն է վնասում}
  \begin{columns}[c]
    \begin{column}{0.52\textwidth}
      Մշտական, քրոնիկ տագնապը (anxiety disorder) ազդում է նաև մարմնի վրա.\par\smallskip
      \begin{itemize}
        \item<2-> \emoji{anatomicalheart}\ բարձր զարկերակային ճնշում
        \item<2-> \emoji{muscle}\ լարվածություն, \emoji{cold}\ դող
        \item<2-> \emoji{forkknife}\ քաշի կորուստ կամ ավելացում
      \end{itemize}
      \onslide<3->{Անվնաս բաներով տարվելով՝ ուղեղը մեզ \textbf{իրական} վնաս է տալիս։}
    \end{column}
    \begin{column}{0.42\textwidth}
      \onslide<4->{%
      \begin{bookquote}
        Ո՞ւմ են պետք գիշատիչներ, երբ ունենք մեր մեծ գանգերը, որոնք մշտական սթրեսով մեզ ներքև են քաշում։
      \end{bookquote}
      \centering\mascot[1.8cm]{sad}}
    \end{column}
  \end{columns}
\end{frame}

% How common: the book's UK numbers, plus WHO's global figure
\begin{frame}{Որքա՞ն տարածված է}
  \begin{columns}[c]
    \begin{column}{0.5\textwidth}
      \centering
      \begin{tikzpicture}
        \foreach \i in {0,...,9}{
          \ifnum\i=0 \def\pc{amOrange}\else \def\pc{plum!22!cream}\fi
          \fill[\pc] (\i*0.62,0.62) circle (0.17);
          \fill[\pc, rounded corners=1.3mm] (\i*0.62-0.22,0) rectangle (\i*0.62+0.22,0.42);}
      \end{tikzpicture}\par\medskip
      \raggedright
      \begin{itemize}
        \item Մեծ Բրիտանիայում \textbf{10 մեծահասակից մոտ 1-ը} կյանքում գոնե մեկ անգամ տագնապային խանգարում է ունենում
        \item<2-> 1993-ից 2007 թվականներին դրանք շատացել են \textbf{12.8\%-ով}. մոտ մեկ միլիոն մարդ ավելի
      \end{itemize}
    \end{column}
    \begin{column}{0.44\textwidth}
      \onslide<3->{%
      \begin{scienceupdate}
        Ամբողջ աշխարհում 2019-ին տագնապային խանգարում ուներ \textbf{301 միլիոն} մարդ։ Դրանք ամենատարածված հոգեկան խանգարումներն են։
      \end{scienceupdate}}
    \end{column}
  \end{columns}
  \only<3->{\source{ԱՀԿ, Anxiety disorders, 2023 (2019-ի տվյալներ)}}
\end{frame}

% ---------------------------------------------------------------- section 1: superstition and conspiracies
\partaccent{lavender}
\topicopener{clover}{Սնահավատություն}{Ի՞նչ ընդհանուր բան ունեն քառատերև երեքնուկն ու ՉԹՕ-ները}{Սնահավատություն, դավադրության տեսություններ ու այլ տարօրինակ համոզմունքներ}{confused}

\begin{frame}{Մի քիչ սնահավատ ենք բոլորս}
  \begin{columns}[c]
    \begin{column}{0.5\textwidth}
      \begin{poll}
        Բարձրացրու ձեռքդ, եթե գոնե մեկն արել ես.\par\smallskip
        \emoji{wood}\ թակել ես փայտին, որ «աչք չդիպչի»\par
        \emoji{ladder}\ չես անցել սանդուղքի տակով\par
        \emoji{blackcat}\ շրջանցել ես սև կատվին\par
        \emoji{calendar}\ ուրբաթ, 13-ին տանից դուրս չես եկել
      \end{poll}
    \end{column}
    \begin{column}{0.44\textwidth}
      \onslide<2->{%
      \begin{card}[colback=lavender!35!cardwhite]
        {\blackfont Սնահավատության նշանը}\par\smallskip
        Լարվում ես իրավիճակներից, որոնք իրականում \textbf{հիմք չունեն}, ու անում ես բաներ, որոնք \textbf{չեն կարող} ազդել դեպքերի վրա։ Միայն թե ավելի ապահով զգաս։
      \end{card}}
    \end{column}
  \end{columns}
\end{frame}

% Burnett's own "involvement" in conspiracies
\begin{frame}{Խոստովանություն. Բըրնեթը դավադիր է}
  \begin{columns}[c]
    \begin{column}{0.62\textwidth}
      {\blackfont Իմ մասին հետաքրքիր փաստ}. ես բազմաթիվ գաղտնի դավադրությունների մասնակից եմ.\par\smallskip
      \begin{itemize}
        \item<2-> \emoji{pill}\ «Մեծ ֆարմայի» հետ թաքցնում եմ քաղցկեղի բուժումները՝ հանուն շահույթի
        \item<3-> \emoji{rocket}\ Թաքցնում եմ, որ Լուսնի վայրէջքը բեմադրված էր
        \item<4-> \emoji{brain}\ Հոգեբուժությունը (իմ աշխատանքը) խարդախություն է՝ ազատ մտածողներին ճնշելու համար
        \item<5-> \emoji{earth}\ Էվոլյուցիան, պատվաստումներն ու գնդաձև Երկիրը գիտնականների «առասպելներն» են
      \end{itemize}
    \end{column}
    \begin{column}{0.32\textwidth}
      \centering
      \onslide<6->{%
      \begin{bookquote}
        Ի վերջո, Երկրի վրա ոչ ոք ավելի հարուստ ու հզոր չէ, քան գիտնականները։
      \end{bookquote}
      \tikz\node[stamp=pinkdark, font=\blackfont\Large] {ԳԱՂՏՆԻ};}
    \end{column}
  \end{columns}
\end{frame}

% How he found out: the comment section
\begin{frame}{Մեկնաբանությունների բաժինն ամեն ինչ գիտի}
  \begin{columns}[c]
    \begin{column}{0.52\textwidth}
      \begin{tikzpicture}
        \node[chatmsg=6.2cm, fill=cardwhite] (c1) at (0,0)
          {\textbf{@ճշմարտություն\_777}\par Աշխարհի ամենավատ գրողն է։ Ու ֆարմայի գործակալ։};
        \node[chatmsg=6.2cm, fill=lavender!35!cardwhite, below=2.5mm of c1, visible on=<2->] (c2)
          {\textbf{@արթնացեք}\par Հակատրանսգենդեր դավադրության մաս է. նրանց բավականաչափ կտրուկ չպաշտպանեց։};
        \node[chatmsg=6.2cm, fill=peach!40!cardwhite, below=2.5mm of c2, visible on=<3->] (c3)
          {\textbf{@ոչխար\_չեմ}\par Պրոտրանսգենդեր դավադրության մաս է. նրանց պաշտպանեց։};
      \end{tikzpicture}
    \end{column}
    \begin{column}{0.42\textwidth}
      Բըրնեթն իր «մասնակցության» մասին իմացավ Guardian-ում իր հոդվածների տակի մեկնաբանություններից։\par\medskip
      \onslide<4->{%
      \begin{bookquote}
        Ոչ միայն մասնակցում եմ բազմաթիվ դավադրությունների, այլև այդ ընթացքում ակտիվորեն պայքարում եմ ինքս իմ դեմ։
      \end{bookquote}}
    \end{column}
  \end{columns}
\end{frame}

% Meme: who is really behind it
\begin{frame}{Իսկ ո՞վ է իրականում կանգնած դրա հետևում}
  \begin{columns}[c]
    \begin{column}{0.48\textwidth}
      Եթե հոդվածը քննադատում է քո հայացքները, հեշտ է եզրակացնել, որ դրա հետևում ճշմարտությունը լռեցնող չար ուժ է։\par\medskip
      \onslide<2->{%
      \begin{card}[colback=lavender!35!cardwhite]
        Իրականում՝ մի տղա է, որ վաղաժամ ճաղատանում է ու նստած է բազմոցին Քարդիֆում։
      \end{card}}
    \end{column}
    \begin{column}{0.46\textwidth}
      \centering
      \onslide<2->{%
      \begin{meme}[4.5cm]{scooby_doo_mask_reveal.jpg}
        \memecap{0.27}{0.6}{ստվերային\\ համաշխարհային\\ իշխանություն}
        \memecap{0.3}{0.1}{բազմոցին նստած\\ ճաղատացող տղա\\ Քարդիֆից}
      \end{meme}}
    \end{column}
  \end{columns}
\end{frame}

% The internet, and the question of the section title
\begin{frame}{Ինտերնետը՝ դավադրությունների պարարտ հող}
  \begin{columns}[c]
    \begin{column}{0.52\textwidth}
      \begin{itemize}
        \item Այսօր կարելի է, տնից դուրս չգալով, գտնել «ապացույցներ» 9/11-ի կամ «ԿՀՎ-ն ու ՁԻԱՀ-ը» թեմայով ու կիսվել համախոհների հետ
        \item<2-> Բայց դավադրության տեսությունները \textbf{նոր չեն}. դրանք միշտ են եղել
        \item<3-> Ուրեմն գուցե ուղեղի ինչ-որ առանձնահատկությունն է մեզ այդքան հակված դարձնում պարանոյիկ երևակայությունների
      \end{itemize}
    \end{column}
    \begin{column}{0.42\textwidth}
      \onslide<4->{%
      \begin{card}[colback=lavender!35!cardwhite]
        \centering
        \bigemoji[1.1cm]{ufo}\quad{\blackfont\Large ?}\quad\bigemoji[1.1cm]{clover}\par\smallskip
        Ներխուժել «Զոնա 51», որ գտնես այլմոլորակայիններին, թե՞ հավատալ քառատերև երեքնուկի հաջողությանը։ Ի՞նչ ընդհանուր բան ունեն։
      \end{card}}
    \end{column}
  \end{columns}
  \source{Aaronovitch \& Langton, Voodoo Histories, 2010}
\end{frame}

% Apophenia, with the book's two examples
\begin{frame}{Ապոֆենիա. կապեր, որոնք չկան}
  \begin{card}[colback=lavender!35!cardwhite]
    \centering \textbf{Ապոֆենիա (apophenia)}՝ կապեր տեսնել այնտեղ, որտեղ դրանք իրականում չկան
  \end{card}
  \vspace{1mm}
  \centering
  \begin{tikzpicture}[ex/.style={pbox, text width=3.6cm, minimum height=1.1cm}, ar/.style={-{Stealth[length=2.6mm]}, plum, line width=1.2pt}]
    \node[ex, visible on=<2->] (a1) at (0,0) {\emoji{jeans}\ Պատահաբար ներքնաշորդ շրջած ես հագել};
    \node[ex, visible on=<2->] (a2) at (4.4,0) {\emoji{moneybag}\ Նույն օրը վիճակախաղում շահեցիր};
    \node[ex, visible on=<2->, fill=lavender!40!cardwhite] (a3) at (8.8,0) {Այսուհետ տոմս ես գնում միայն \textbf{շրջած} ներքնաշորով};
    \draw[ar, visible on=<2->] (a1) -- (a2); \draw[ar, visible on=<2->] (a2) -- (a3);
    \node[ex, visible on=<3->] (b1) at (0,-1.35) {\emoji{skull}\ Երկու հայտնի մարդ մեկ ամսում մահացավ};
    \node[ex, visible on=<3->] (b2) at (4.4,-1.35) {\emoji{newspaper}\ Երկուսն էլ քննադատել էին կառավարությանը};
    \node[ex, visible on=<3->, fill=lavender!40!cardwhite] (b3) at (8.8,-1.35) {«Նրանց \textbf{սպանել} են»};
    \draw[ar, visible on=<3->] (b1) -- (b2); \draw[ar, visible on=<3->] (b2) -- (b3);
  \end{tikzpicture}\par\smallskip
  \raggedright
  \onslide<4->{Ամեն սնահավատության կամ դավադրության հիմքում, հավանաբար, \textbf{անկապ} դեպքերի միջև «գտնված» կապ է։ Ու դա կարող է պատահել ամեն մեկի հետ։}
  \source{Fyfe et al., Cortex, 2008}
\end{frame}

% Try it: the Face on Mars (picture added, not in the book)
\begin{frame}{Փորձիր. ի՞նչ ես տեսնում}
  \begin{columns}[T]
    \begin{column}{0.34\textwidth}
      \centering
      \includegraphics[height=2.7cm]{mars_face_viking_1976.jpg}\par\smallskip
      {\small Մարս, 1976, Viking 1}\par\smallskip
      \begin{tryit}
        Գրիր չաթում՝ ի՞նչ ես տեսնում։
      \end{tryit}
    \end{column}
    \begin{column}{0.3\textwidth}
      \centering
      \onslide<2->{%
      \includegraphics[height=2.9cm]{mars_face_hirise_2007.jpg}\par\smallskip
      {\small Նույն բլուրը, 2007, շատ ավելի լավ տեսախցիկով}}
    \end{column}
    \begin{column}{0.3\textwidth}
      \onslide<3->{%
      \begin{card}[colback=lavender!35!cardwhite]
        Ուղեղը դեմք «գտավ» լույսի ու ստվերի խաղի մեջ։ Ոմանք տարիներով պնդում էին, որ այն այլմոլորակայիններն են կառուցել։
      \end{card}}
    \end{column}
  \end{columns}
  \source{NASA/JPL: Viking 1, 1976; MRO HiRISE, 2007}
\end{frame}

% Why the brain looks for patterns: too much information, too little time
\begin{frame}{Ուղեղը՝ տեղեկատվության հեղեղի տակ}
  \begin{columns}[c]
    \begin{column}{0.55\textwidth}
      \begin{itemize}
        \item Ամեն պահ ուղեղ է հասնում չափազանց խիտ տեղեկատվություն, որի միայն մի մասն է կարևոր
        \item<2-> Մշակելու համար ունի \textbf{վայրկյանի մի մասը}
        \item<3-> Դրա համար ուղեղն օգտագործում է \textbf{կարճ ճանապարհներ} (shortcuts). օրինակ՝ \textbf{օրինաչափություններ} փնտրել
      \end{itemize}
    \end{column}
    \begin{column}{0.4\textwidth}
      \centering
      \begin{tikzpicture}[ar/.style={-{Stealth[length=2.2mm]}, plum, line width=1.1pt}]
        \node[pbox=cardwhite, text width=4.2cm] (r) at (0,0) {\emoji{eye}\ ցանցաթաղանթ (retina)};
        \node[pbox=cardwhite, text width=4.2cm] (v) at (0,-0.9) {տեսողական կեղև};
        \node[pbox=cardwhite, text width=4.2cm] (h) at (0,-1.8) {հիպոկամպ};
        \node[pbox=cardwhite, text width=4.2cm] (p) at (0,-2.7) {նախաճակատային կեղև};
        \draw[ar] (r) -- (v); \draw[ar] (v) -- (h); \draw[ar] (h) -- (p);
        \node[tag, text width=4.4cm, font=\small] at (0,-3.45) {շատ գոտիներ՝ \textbf{միասին}};
      \end{tikzpicture}
    \end{column}
  \end{columns}
  \vspace{1mm}
  \onslide<4->{%
  \begin{card}[colback=lavender!35!cardwhite]
    \emoji{newspaper}\ «Գտնվել է X-ի գոտին» տիպի լուրերը մոլորեցնող են. լավագույն դեպքում դա կիսատ բացատրություն է։
  \end{card}}
\end{frame}

% False alarms are cheap; but they add up
\begin{frame}{Կեղծ ահազանգերի գինը}
  \begin{bookquote}
    Ի՞նչ արժեն մի քանի կեղծ ահազանգը, եթե դրանով գոյատևումդ ապահովված է։
  \end{bookquote}
  \vspace{1mm}
  \centering
  \begin{tikzpicture}[ar/.style={font=\blackfont\LARGE, text=plum}]
    \node[pbox=lavender!40!cardwhite, text width=2.7cm, minimum height=1.3cm, visible on=<2->] (a) at (0,0) {\emoji{eyes}\\ ապոֆենիա};
    \node[ar, visible on=<3->] at (1.85,0) {+};
    \node[pbox=brainpink!35!cardwhite, text width=2.7cm, minimum height=1.3cm, visible on=<3->] (b) at (3.7,0) {\emoji{runner}\\ «կռվիր կամ փախիր»};
    \node[ar, visible on=<4->] at (5.55,0) {+};
    \node[pbox=peach!40!cardwhite, text width=2.7cm, minimum height=1.3cm, visible on=<4->] (c) at (7.4,0) {\emoji{skull}\\ ամենավատ սցենարը};
    \node[ar, visible on=<5->] at (9.25,0) {=};
    \node[pbox=butter!60!cardwhite, text width=2.2cm, minimum height=1.3cm, visible on=<5->] (d) at (10.8,0) {\emoji{scream}\\ գլուխդ լիքն է};
  \end{tikzpicture}\par\medskip
  \raggedright
  \onslide<6->{Ահա թե ինչու են գրեթե բոլոր սնահավատությունները \textbf{դժբախտությունից խուսափելու} մասին։}
\end{frame}

% No conspiracy ever helps anyone
\begin{frame}{Բարի դավադրություններ չեն լինում}
  \begin{columns}[c]
    \begin{column}{0.38\textwidth}
      \begin{bookquote}
        Երբեք չես լսի դավադրության մասին, որի նպատակը մարդկանց օգնելն է։ Խորհրդավոր էլիտան բարեգործական տոնավաճառներ չի կազմակերպում։
      \end{bookquote}
    \end{column}
    \begin{column}{0.58\textwidth}
      \centering
      \begin{tikzpicture}
        \draw[plum, line width=2pt] (-2.3,-1.2) -- (-2.3,-2.2) (2.3,-1.2) -- (2.3,-2.2);
        \fill[plum, rounded corners=2mm] (-3.4,-1.35) rectangle (3.6,1.45);
        \filldraw[fill=butter!70!cardwhite, draw=plum, line width=1pt, rounded corners=2mm] (-3.5,-1.25) rectangle (3.5,1.55);
        \node[font=\blackfont\normalsize, text=plum, align=center] at (0,0.85) {Իլյումինատների\\ բարեգործական տոնավաճառ};
        \node[font=\small, text=plum, align=center] at (0,-0.45) {\emoji{cake}\ Տնական թխվածք՝ 500 դրամ\\ Ամբողջ հասույթը՝ որբանոցին};
        \node[stamp=pinkdark, font=\blackfont\LARGE, visible on=<2->] at (1.6,-1.45) {ՉԻ ԼԻՆՈՒՄ};
      \end{tikzpicture}
    \end{column}
  \end{columns}
\end{frame}

% Childhood: parents' superstitions and a world run by powerful figures
\begin{frame}{Ամեն ինչ սկսվում է մանկությունից}
  \begin{columns}[c]
    \begin{column}{0.58\textwidth}
      \begin{itemize}
        \item Մանկությունում գրեթե ամեն ինչ իմանում ենք մեծերից, ու ծնողների սնահավատություններն անցնում են մեզ առանց որևէ «ապացույցի»
        \item<2-> Աշխարհը թվում է \textbf{կազմակերպված}. այն ղեկավարում են հզոր, անհասկանալի մարդիկ
        \item<3-> Մեծանալով՝ ոմանց համար հանգստացնող է մտածել, որ հիմա էլ է այդպես։ «Հզորները» մագնատներն են, մարդակեր այլմոլորակային մողեսները կամ գիտնականները
      \end{itemize}
    \end{column}
    \begin{column}{0.36\textwidth}
      \centering
      \bigemoji[1.1cm]{family}\ \bigemoji[1.1cm]{lizard}\par\medskip
      \onslide<4->{%
      \begin{chatbox}
        Ո՞ր սնահավատությունն ես «ժառանգել» ընտանիքից։
      \end{chatbox}}
    \end{column}
  \end{columns}
  \source{Leonard, 1989}
\end{frame}

% The brain cannot stand randomness
\begin{frame}{Ուղեղը չի հանդուրժում պատահականությունը}
  \begin{columns}[c]
    \begin{column}{0.56\textwidth}
      \begin{itemize}
        \item Ուղեղի համար դժվար է ընդունել, որ ինչ-որ բան կարող է պատահել \textbf{առանց պատճառի}. այդ դեպքում ոչինչ չես կարող անել
        \item<2-> Այստեղից՝ «ճակատագիրը»։ Անհաջողությունների շա՞րք. երևի կոտրված հայելու մեղքն է \emoji{mirror}
        \item<3-> Կամ չարաճճի փերիներն են. երկաթ չեն սիրում, ուրեմն պայտ կախիր
      \end{itemize}
    \end{column}
    \begin{column}{0.38\textwidth}
      \centering
      \onslide<4->{%
      \bigemoji[1.0cm]{airplane}\par\smallskip
      \begin{bookquote}
        Ավելի լավ է՝ ղեկին հարբած օդաչու լինի, քան ընդհանրապես ոչ ոք։
      \end{bookquote}}
    \end{column}
  \end{columns}
  \vspace{2mm}
  \onslide<5->{\centering\small Պատահականությամբ «գնացող» աշխարհն ավելի վախեցնող է, քան գաղտնի էլիտայով աշխարհը։\par}
\end{frame}

% Locus of control: rituals vs campaigns
\begin{frame}{Վերահսկողության լոկուս}
  \begin{card}[colback=lavender!35!cardwhite]
    \textbf{Վերահսկողության լոկուս (locus of control)}՝ որքանով ես \textbf{հավատում}, որ դու ես կառավարում քեզ հետ կատարվողը (իրականում որքան՝ կարևոր չէ)
  \end{card}
  \vspace{1mm}
  \begin{columns}[T]
    \begin{column}{0.47\textwidth}
      \onslide<2->{%
      \begin{card}
        \emoji[1.3em]{clover}\ {\blackfont Սնահավատը}\par
        Աղ է շպրտում, թակում փայտին ու հանգստանում. աղետը կանխված է։ Շարունակում է իր օրը։
      \end{card}}
    \end{column}
    \begin{column}{0.47\textwidth}
      \onslide<3->{%
      \begin{card}
        \emoji[1.3em]{ufo}\ {\blackfont Դավադրության տեսաբանը}\par
        «Արթնացնում» է ուրիշներին, փորփրում է «ավելի խորը», չհամաձայնողները՝ «անմիտ ոչխարներ» \emoji{sheep}
      \end{card}}
    \end{column}
  \end{columns}
  \onslide<4->{\centering\small Ե՞րբ են քեզ վերջին անգամ համոզել նապաստակի թաթի «թաքնված ճշմարտության» մեջ\par}
\end{frame}

% Chicken and egg: hippocampus, cortisol, control
\begin{frame}{Հավը, ձուն ու հիպոկամպը}
  \begin{columns}[c]
    \begin{column}{0.53\textwidth}
      \centering
      \begin{tikzpicture}[ar/.style={-{Stealth[length=2.4mm]}, plum, line width=1.1pt},
          hb/.style={pbox, text width=2.8cm, font=\footnotesize}]
        \node[hb, fill=lavender!40!cardwhite] (hip) at (0,0) {\emoji{brain}\ մեծ հիպոկամպ};
        \node[hb, fill=mint!40!cardwhite] (ctl) at (4.1,0) {վերահսկողության զգացում};
        \draw[{Stealth[length=2.4mm]}-{Stealth[length=2.4mm]}, plum, line width=1.1pt] (hip) -- node[above, tag, font=\scriptsize, yshift=1mm] {կապ} (ctl);
        \node[hb, fill=brainpink!35!cardwhite, visible on=<2->] (low) at (4.1,-1.8) {քիչ վերահսկողություն $\rightarrow$ սթրես};
        \node[hb, fill=peach!45!cardwhite, visible on=<2->] (cort) at (0,-1.8) {կորտիզոլը փոքրացնում է հիպոկամպը};
        \draw[ar, visible on=<2->] (low) -- (cort);
        \draw[ar, visible on=<2->] (cort) -- (hip);
      \end{tikzpicture}\par\smallskip
      \onslide<3->{{\blackfont\small Ո՞րն է պատճառը, ո՞րը՝ հետևանքը}\par
      {\small «Ուղեղը երբեք ոչինչ հեշտ չի դարձնում»}}
    \end{column}
    \begin{column}{0.44\textwidth}
      \onslide<4->{%
      \begin{scienceupdate}
        Նոր հետազոտություններում դավադրության տեսություններին ավելի հակված են նրանք, ովքեր \textbf{չեն} զգում, որ վերահսկում են իրենց կյանքը։ Իսկ տեսությունները այդ զգացումը միայն թուլացնում են։
      \end{scienceupdate}}
    \end{column}
  \end{columns}
  \only<1-3>{\source{Pruessner et al., 2005; O'Brien et al., 2004}}%
  \only<4->{\source{Douglas, Sutton \& Cichocka, Curr. Dir. Psychol. Sci., 2017}}
\end{frame}

% Hard to change minds, and hard to study
\begin{frame}{Ինչու՞ է դժվար ուսումնասիրել}
  \begin{columns}[c]
    \begin{column}{0.52\textwidth}
      \begin{itemize}
        \item Խոր համոզմունքը դժվար է փոխել, որքան էլ ապացույց բերես
        \item<2-> Իսկ ի՞նչ է սնահավատությունը հոգեբանորեն. համոզմու՞նք, գաղափա՞ր
        \item<3-> Ինչպե՞ս կերևա այն ուղեղի ակտիվության մեջ
      \end{itemize}
    \end{column}
    \begin{column}{0.42\textwidth}
      \centering
      \onslide<4->{%
      \bigemoji[1.0cm]{musical}\par\smallskip
      \begin{bookquote}
        Ակտիվությունը տեսնելը չի նշանակում հասկանալ, թե ինչ է այն ներկայացնում. ինչպես դաշնամուրի ստեղները տեսնելը չի նշանակում Մոցարտ նվագել։
      \end{bookquote}}
    \end{column}
  \end{columns}
\end{frame}

% Lindeman et al.: design
\begin{frame}{Երկու բալ ու նշան տիեզերքից}
  \begin{columns}[c]
    \begin{column}{0.6\textwidth}
      \begin{itemize}
        \item Լինդեմանն ու գործընկերները fMRI-ով սկանավորեցին \textbf{12} գերբնականին հավատացողի ու \textbf{11} թերահավատի
        \item<2-> Նախ պատկերացնում էին ծանր պահ. աշխատանքի կորուստ, բաժանում
        \item<3-> Հետո նայում էին նկարներ. լեռան գագաթ, \textbf{իրար կապված երկու բալ}
      \end{itemize}
      \onslide<4->{%
      \begin{card}[colback=lavender!35!cardwhite]
        \emoji{crystalball}\ \textbf{Հավատացողը}. «Բալերը կապված են, ուրեմն մեր կապն ամուր է»\par
        \onslide<5->{\emoji{monocle}\ \textbf{Թերահավատը}. «Երկու բալ է»}
      \end{card}}
    \end{column}
    \begin{column}{0.34\textwidth}
      \centering
      \onslide<3->{%
      \begin{tikzpicture}[scale=0.85]
        \draw[mintdark, line width=1.6pt] (-0.75,-0.1) .. controls (-0.6,1.0) and (-0.2,1.6) .. (0.15,2.1);
        \draw[mintdark, line width=1.6pt] (0.85,-0.1) .. controls (0.75,1.0) and (0.4,1.6) .. (0.15,2.1);
        \fill[mintdark, rotate around={35:(0.15,2.1)}] (0.6,2.1) ellipse (0.45 and 0.18);
        \draw[pinkdark, line width=2.4pt] (-0.42,1.1) -- (0.62,1.1);
        \draw[pinkdark, line width=1.8pt] (0.1,1.1) -- (-0.15,0.65) (0.1,1.1) -- (0.35,0.65);
        \filldraw[fill=amRed, draw=plum, line width=0.8pt] (-0.8,-0.55) circle (0.55);
        \filldraw[fill=amRed, draw=plum, line width=0.8pt] (0.9,-0.55) circle (0.55);
        \fill[white, opacity=0.7] (-0.98,-0.35) circle (0.11) (0.72,-0.35) circle (0.11);
      \end{tikzpicture}}
    \end{column}
  \end{columns}
  \source{Lindeman et al., Soc. Cogn. Affect. Neurosci., 2013}
\end{frame}

% Lindeman et al.: results (the book says temporal; the paper says frontal)
\begin{frame}{Թերահավատի արգելակը}
  \begin{columns}[c]
    \begin{column}{0.36\textwidth}
      \centering
      \begin{tikzpicture}[scale=0.85]
        \filldraw[fill=brainpink!35!cardwhite, draw=plum, line width=1pt] (-1.02,0) ellipse (0.98 and 1.7);
        \filldraw[fill=brainpink!35!cardwhite, draw=plum, line width=1pt] (1.02,0) ellipse (0.98 and 1.7);
        \draw[fold, line width=0.8pt] (-1.5,0.9) .. controls (-1.0,0.6) .. (-0.6,1.0) (-1.6,-0.2) .. controls (-1.1,-0.4) .. (-0.5,-0.1)
          (1.5,0.9) .. controls (1.0,0.6) .. (0.6,1.0) (1.6,-0.2) .. controls (1.1,-0.4) .. (0.5,-0.1);
        \node[tag, font=\blackfont\scriptsize] at (0,2.0) {ճակատ};
        \node[tag, font=\small] at (-1.0,-2.0) {ձախ};
        \node[tag, font=\small] at (1.0,-2.0) {աջ};
        \filldraw[fill=amOrange, draw=plum, line width=0.8pt, visible on=<2->] (-1.6,0.95) circle (0.27);
        \filldraw[fill=amBlue, draw=plum, line width=0.8pt, visible on=<3->] (1.6,0.95) circle (0.27);
      \end{tikzpicture}
    \end{column}
    \begin{column}{0.6\textwidth}
      \onslide<2->{%
      \begin{card}
        \textcolor{amOrange}{\Large$\bullet$}\ \textbf{Ձախ ստորին ճակատային գալար} (inferior frontal gyrus)՝ ակտիվ էր \textbf{բոլորի} մոտ
      \end{card}}
      \onslide<3->{%
      \begin{card}
        \textcolor{amBlue}{\Large$\bullet$}\ \textbf{Աջը}՝ ավելի ակտիվ էր \textbf{թերահավատների} մոտ։ Այն կապված է ճանաչողական \textbf{արգելակման} (cognitive inhibition) հետ
      \end{card}}
      \onslide<4->{%
      \begin{scienceupdate}
        Գրքում՝ «քունքային» գալար։ Հոդվածում ճակատային գալարն է։
      \end{scienceupdate}}
    \end{column}
  \end{columns}
  \source{Lindeman et al., 2013; Hampshire et al., NeuroImage, 2010}
\end{frame}

% Limits of the study, and why conspiracy theorists are hard to scan
\begin{frame}{Իսկ դավադրության տեսաբաննե՞րը}
  \begin{columns}[c]
    \begin{column}{0.42\textwidth}
      \begin{itemize}
        \item Ընդամենը 23 մարդ
        \item<2-> Ինչպե՞ս չափել «գերբնականին հակվածությունը». մետրական համակարգում այդպիսի միավոր չկա
        \item<3-> Դավադրության տեսաբաններին ուսումնասիրելն ավելի դժվար է. նրանք գաղտնապահ են ու չեն վստահում պաշտոնական կառույցներին
      \end{itemize}
    \end{column}
    \begin{column}{0.54\textwidth}
      \centering
      \onslide<4->{%
      \begin{tikzpicture}
        \node at (-2.3,-1.7) {\bigemoji[1.2cm]{scientist}};
        \node at (2.4,-1.7) {\bigemoji[1.2cm]{detective}};
        \node[say, text width=4.3cm, font=\mediumfont\small, callout absolute pointer={(-2.2,-1.05)}] at (-0.9,0.75)
          {Եկեք մեր գաղտնի լաբորատորիա։ Կփակենք ձեզ մետաղե խողովակում ու կսկանավորենք ուղեղը։};
        \node[say, font=\handfont\Large, callout absolute pointer={(2.3,-1.1)}, visible on=<5->] at (2.7,0.0) {Ո՛չ};
      \end{tikzpicture}}
    \end{column}
  \end{columns}
\end{frame}

% Closing joke of the section
\begin{frame}{Բայց ես այդպես էլ կասեի, չէ՞}
  \begin{columns}[c]
    \begin{column}{0.56\textwidth}
      Այս բաժնում ներկայացվածը խելամիտ տեսություններ ու ենթադրություններ են՝ հիմնված այսօրվա տվյալների վրա։\par\medskip
      \onslide<2->{%
      \begin{bookquote}
        Բայց ես այդպես էլ կասեի, չէ՞։ Այս ամբողջ գլուխը կարող է լինել մարդկանց խավարի մեջ պահելու դավադրության մի մասը...
      \end{bookquote}}
    \end{column}
    \begin{column}{0.38\textwidth}
      \centering
      \begin{tikzpicture}
        \node {\mascot[3.2cm]{smug}};
        \node[stamp=pinkdark, font=\blackfont\LARGE, visible on=<3->] at (0.1,-0.4) {ԳԱՂՏՆԻ};
      \end{tikzpicture}
    \end{column}
  \end{columns}
\end{frame}
